\documentclass[11pt,a4paper]{article}
\usepackage[utf8]{inputenc}
\usepackage[margin=1in]{geometry}
\usepackage{amsmath,amssymb,amsthm}
\usepackage{graphicx}
\usepackage{hyperref}
\usepackage{enumitem}
\usepackage{tcolorbox}
\usepackage{fancyhdr}
\usepackage{algorithm}
\usepackage{algpseudocode}
\usepackage{listings}
\usepackage{xcolor}

% Page style
\pagestyle{fancy}
\fancyhf{}
\rhead{GATE CSE Study Guide}
\lhead{\leftmark}
\cfoot{\thepage}

% Custom environments
\newtcolorbox{defbox}{colback=blue!5!white,colframe=blue!75!black,title=Definition}
\newtcolorbox{thmbox}{colback=green!5!white,colframe=green!75!black,title=Theorem}
\newtcolorbox{notebox}{colback=yellow!5!white,colframe=yellow!75!black,title=Important Note}
\newtcolorbox{formulabox}{colback=red!5!white,colframe=red!75!black,title=Formula}

\title{\textbf{Comprehensive Study Guide for GATE CSE}\\
\large All Theoretical Concepts}
\author{Prepared for GATE Computer Science Engineering}
\date{\today}

\begin{document}

\maketitle
\newpage

\tableofcontents
\newpage

%==============================================================================
\section{Engineering Mathematics}
%==============================================================================

\subsection{Discrete Mathematics}

\subsubsection{Mathematical Logic}

\begin{defbox}
\textbf{Proposition:} A declarative statement that is either true or false, but not both.
\end{defbox}

\textbf{Logical Connectives:}
\begin{itemize}
    \item Negation ($\neg$): NOT
    \item Conjunction ($\wedge$): AND
    \item Disjunction ($\vee$): OR
    \item Implication ($\rightarrow$): IF-THEN
    \item Biconditional ($\leftrightarrow$): IF AND ONLY IF
\end{itemize}

\begin{formulabox}
\textbf{De Morgan's Laws:}
\begin{align*}
    \neg(p \wedge q) &\equiv \neg p \vee \neg q \\
    \neg(p \vee q) &\equiv \neg p \wedge \neg q
\end{align*}
\end{formulabox}

\textbf{Tautology:} A proposition that is always true.

\textbf{Contradiction:} A proposition that is always false.

\textbf{Contingency:} A proposition that is neither a tautology nor a contradiction.

\subsubsection{Set Theory}

\begin{defbox}
\textbf{Set:} A well-defined collection of distinct objects.
\end{defbox}

\textbf{Set Operations:}
\begin{itemize}
    \item Union: $A \cup B = \{x : x \in A \text{ or } x \in B\}$
    \item Intersection: $A \cap B = \{x : x \in A \text{ and } x \in B\}$
    \item Difference: $A - B = \{x : x \in A \text{ and } x \notin B\}$
    \item Complement: $A^c = U - A$ (where U is universal set)
    \item Cartesian Product: $A \times B = \{(a,b) : a \in A, b \in B\}$
\end{itemize}

\begin{formulabox}
\textbf{Set Identities:}
\begin{align*}
    A \cup (B \cap C) &= (A \cup B) \cap (A \cup C) \quad \text{(Distributive)} \\
    A \cap (B \cup C) &= (A \cap B) \cup (A \cap C) \quad \text{(Distributive)} \\
    (A \cup B)^c &= A^c \cap B^c \quad \text{(De Morgan's)} \\
    (A \cap B)^c &= A^c \cup B^c \quad \text{(De Morgan's)}
\end{align*}
\end{formulabox}

\textbf{Power Set:} $P(A) = \{S : S \subseteq A\}$. If $|A| = n$, then $|P(A)| = 2^n$.

\subsubsection{Relations}

\begin{defbox}
\textbf{Relation:} A subset of the Cartesian product $A \times B$.
\end{defbox}

\textbf{Properties of Relations (on set A):}
\begin{itemize}
    \item \textbf{Reflexive:} $\forall a \in A, (a,a) \in R$
    \item \textbf{Symmetric:} $(a,b) \in R \Rightarrow (b,a) \in R$
    \item \textbf{Antisymmetric:} $(a,b) \in R \wedge (b,a) \in R \Rightarrow a = b$
    \item \textbf{Transitive:} $(a,b) \in R \wedge (b,c) \in R \Rightarrow (a,c) \in R$
\end{itemize}

\begin{defbox}
\textbf{Equivalence Relation:} A relation that is reflexive, symmetric, and transitive.
\end{defbox}

\begin{defbox}
\textbf{Partial Order:} A relation that is reflexive, antisymmetric, and transitive.
\end{defbox}

\subsubsection{Functions}

\begin{defbox}
\textbf{Function:} A relation $f: A \rightarrow B$ where each element in A is related to exactly one element in B.
\end{defbox}

\textbf{Types of Functions:}
\begin{itemize}
    \item \textbf{Injective (One-to-One):} $f(a_1) = f(a_2) \Rightarrow a_1 = a_2$
    \item \textbf{Surjective (Onto):} $\forall b \in B, \exists a \in A$ such that $f(a) = b$
    \item \textbf{Bijective:} Both injective and surjective
\end{itemize}

\subsubsection{Combinatorics}

\begin{formulabox}
\textbf{Permutations:}
\begin{align*}
    P(n,r) &= \frac{n!}{(n-r)!} \\
    \text{Circular Permutations} &= (n-1)!
\end{align*}
\end{formulabox}

\begin{formulabox}
\textbf{Combinations:}
\[
    C(n,r) = \binom{n}{r} = \frac{n!}{r!(n-r)!}
\]
\end{formulabox}

\begin{formulabox}
\textbf{Binomial Theorem:}
\[
    (x+y)^n = \sum_{k=0}^{n} \binom{n}{k} x^{n-k} y^k
\]
\end{formulabox}

\textbf{Pigeonhole Principle:} If $n$ items are placed into $m$ containers with $n > m$, then at least one container must contain more than one item.

\textbf{Inclusion-Exclusion Principle:}
\[
    |A \cup B| = |A| + |B| - |A \cap B|
\]

\subsubsection{Graph Theory}

\begin{defbox}
\textbf{Graph:} $G = (V, E)$ where V is a set of vertices and E is a set of edges.
\end{defbox}

\textbf{Graph Properties:}
\begin{itemize}
    \item \textbf{Degree:} Number of edges incident to a vertex
    \item \textbf{Handshaking Lemma:} $\sum_{v \in V} deg(v) = 2|E|$
    \item \textbf{Path:} Sequence of vertices where each consecutive pair is connected by an edge
    \item \textbf{Cycle:} A path that starts and ends at the same vertex
    \item \textbf{Connected Graph:} There exists a path between every pair of vertices
\end{itemize}

\textbf{Special Graphs:}
\begin{itemize}
    \item \textbf{Complete Graph $K_n$:} Every pair of vertices is connected
    \item \textbf{Bipartite Graph:} Vertices can be divided into two disjoint sets
    \item \textbf{Tree:} Connected acyclic graph with $|E| = |V| - 1$
\end{itemize}

\begin{thmbox}
\textbf{Euler's Theorem:} A connected graph has an Eulerian circuit if and only if every vertex has even degree.
\end{thmbox}

\begin{thmbox}
\textbf{Hamiltonian Path:} A path that visits every vertex exactly once. (No efficient algorithm exists to find it - NP-complete)
\end{thmbox}

\subsection{Linear Algebra}

\subsubsection{Matrices}

\begin{defbox}
\textbf{Matrix:} A rectangular array of numbers arranged in rows and columns.
\end{defbox}

\textbf{Matrix Operations:}
\begin{itemize}
    \item Addition: $(A + B)_{ij} = A_{ij} + B_{ij}$
    \item Scalar Multiplication: $(cA)_{ij} = c \cdot A_{ij}$
    \item Matrix Multiplication: $(AB)_{ij} = \sum_{k} A_{ik}B_{kj}$
    \item Transpose: $(A^T)_{ij} = A_{ji}$
\end{itemize}

\begin{defbox}
\textbf{Determinant:} A scalar value that can be computed from a square matrix. For $2 \times 2$ matrix:
\[
    \det\begin{pmatrix} a & b \\ c & d \end{pmatrix} = ad - bc
\]
\end{defbox}

\textbf{Properties of Determinants:}
\begin{itemize}
    \item $\det(AB) = \det(A) \cdot \det(B)$
    \item $\det(A^T) = \det(A)$
    \item If a row/column is all zeros, $\det(A) = 0$
    \item $\det(A^{-1}) = \frac{1}{\det(A)}$
\end{itemize}

\subsubsection{Vector Spaces}

\begin{defbox}
\textbf{Vector Space:} A set V with operations of addition and scalar multiplication satisfying 8 axioms (closure, associativity, commutativity, identity, inverse, distributivity).
\end{defbox}

\begin{defbox}
\textbf{Linear Independence:} Vectors $v_1, v_2, \ldots, v_n$ are linearly independent if:
\[
    c_1v_1 + c_2v_2 + \cdots + c_nv_n = 0 \Rightarrow c_1 = c_2 = \cdots = c_n = 0
\]
\end{defbox}

\begin{defbox}
\textbf{Basis:} A linearly independent set of vectors that spans the entire vector space.
\end{defbox}

\begin{defbox}
\textbf{Dimension:} The number of vectors in a basis.
\end{defbox}

\subsubsection{Eigenvalues and Eigenvectors}

\begin{defbox}
\textbf{Eigenvalue and Eigenvector:} For matrix A, $\lambda$ is an eigenvalue and $v$ is an eigenvector if:
\[
    Av = \lambda v
\]
\end{defbox}

\textbf{Characteristic Equation:}
\[
    \det(A - \lambda I) = 0
\]

\begin{notebox}
Sum of eigenvalues = Trace of matrix\\
Product of eigenvalues = Determinant of matrix
\end{notebox}

\subsection{Calculus}

\subsubsection{Limits and Continuity}

\begin{defbox}
\textbf{Limit:} $\lim_{x \to a} f(x) = L$ if $f(x)$ approaches L as x approaches a.
\end{defbox}

\textbf{Important Limits:}
\begin{align*}
    \lim_{x \to 0} \frac{\sin x}{x} &= 1 \\
    \lim_{x \to \infty} \left(1 + \frac{1}{x}\right)^x &= e
\end{align*}

\begin{defbox}
\textbf{Continuity:} Function f is continuous at x = a if:
\begin{enumerate}
    \item $f(a)$ is defined
    \item $\lim_{x \to a} f(x)$ exists
    \item $\lim_{x \to a} f(x) = f(a)$
\end{enumerate}
\end{defbox}

\subsubsection{Differentiation}

\begin{formulabox}
\textbf{Basic Derivatives:}
\begin{align*}
    \frac{d}{dx}(x^n) &= nx^{n-1} \\
    \frac{d}{dx}(e^x) &= e^x \\
    \frac{d}{dx}(\ln x) &= \frac{1}{x} \\
    \frac{d}{dx}(\sin x) &= \cos x \\
    \frac{d}{dx}(\cos x) &= -\sin x
\end{align*}
\end{formulabox}

\textbf{Rules of Differentiation:}
\begin{itemize}
    \item Product Rule: $(uv)' = u'v + uv'$
    \item Quotient Rule: $\left(\frac{u}{v}\right)' = \frac{u'v - uv'}{v^2}$
    \item Chain Rule: $(f(g(x)))' = f'(g(x)) \cdot g'(x)$
\end{itemize}

\subsubsection{Integration}

\begin{formulabox}
\textbf{Basic Integrals:}
\begin{align*}
    \int x^n \, dx &= \frac{x^{n+1}}{n+1} + C \quad (n \neq -1) \\
    \int \frac{1}{x} \, dx &= \ln|x| + C \\
    \int e^x \, dx &= e^x + C \\
    \int \sin x \, dx &= -\cos x + C \\
    \int \cos x \, dx &= \sin x + C
\end{align*}
\end{formulabox}

\textbf{Methods of Integration:}
\begin{itemize}
    \item Substitution
    \item Integration by Parts: $\int u \, dv = uv - \int v \, du$
    \item Partial Fractions
\end{itemize}

\subsection{Probability and Statistics}

\subsubsection{Probability Basics}

\begin{defbox}
\textbf{Probability:} $P(A) = \frac{\text{Number of favorable outcomes}}{\text{Total number of outcomes}}$
\end{defbox}

\textbf{Axioms of Probability:}
\begin{enumerate}
    \item $0 \leq P(A) \leq 1$
    \item $P(S) = 1$ (where S is sample space)
    \item If A and B are mutually exclusive: $P(A \cup B) = P(A) + P(B)$
\end{enumerate}

\begin{formulabox}
\textbf{Addition Rule:}
\[
    P(A \cup B) = P(A) + P(B) - P(A \cap B)
\]
\end{formulabox}

\begin{formulabox}
\textbf{Conditional Probability:}
\[
    P(A|B) = \frac{P(A \cap B)}{P(B)}
\]
\end{formulabox}

\begin{formulabox}
\textbf{Bayes' Theorem:}
\[
    P(A|B) = \frac{P(B|A) \cdot P(A)}{P(B)}
\]
\end{formulabox}

\begin{defbox}
\textbf{Independent Events:} Events A and B are independent if $P(A \cap B) = P(A) \cdot P(B)$
\end{defbox}

\subsubsection{Random Variables}

\begin{defbox}
\textbf{Random Variable:} A function that assigns a numerical value to each outcome in a sample space.
\end{defbox}

\begin{defbox}
\textbf{Expected Value (Mean):}
\[
    E[X] = \sum_{i} x_i P(X = x_i) \quad \text{(discrete)}
\]
\end{defbox}

\begin{defbox}
\textbf{Variance:}
\[
    Var(X) = E[(X - \mu)^2] = E[X^2] - (E[X])^2
\]
\end{defbox}

\textbf{Properties:}
\begin{align*}
    E[aX + b] &= aE[X] + b \\
    Var(aX + b) &= a^2 Var(X)
\end{align*}

\subsubsection{Distributions}

\textbf{Binomial Distribution:} $X \sim B(n,p)$
\[
    P(X = k) = \binom{n}{k} p^k (1-p)^{n-k}
\]
Mean: $np$, Variance: $np(1-p)$

\textbf{Poisson Distribution:} $X \sim Poisson(\lambda)$
\[
    P(X = k) = \frac{e^{-\lambda} \lambda^k}{k!}
\]
Mean: $\lambda$, Variance: $\lambda$

\textbf{Normal Distribution:} $X \sim N(\mu, \sigma^2)$
\[
    f(x) = \frac{1}{\sigma\sqrt{2\pi}} e^{-\frac{(x-\mu)^2}{2\sigma^2}}
\]

%==============================================================================
\section{Digital Logic}

\subsection{Boolean Algebra}

\begin{defbox}
\textbf{Boolean Algebra:} An algebraic structure with two binary operations (AND, OR) and one unary operation (NOT).
\end{defbox}

\textbf{Boolean Laws:}
\begin{itemize}
    \item Identity: $A + 0 = A$, $A \cdot 1 = A$
    \item Null: $A + 1 = 1$, $A \cdot 0 = 0$
    \item Idempotent: $A + A = A$, $A \cdot A = A$
    \item Complement: $A + A' = 1$, $A \cdot A' = 0$
    \item Involution: $(A')' = A$
    \item Commutative: $A + B = B + A$, $A \cdot B = B \cdot A$
    \item Associative: $(A + B) + C = A + (B + C)$
    \item Distributive: $A + (B \cdot C) = (A + B) \cdot (A + C)$
    \item Absorption: $A + (A \cdot B) = A$, $A \cdot (A + B) = A$
    \item De Morgan's: $(A + B)' = A' \cdot B'$, $(A \cdot B)' = A' + B'$
\end{itemize}

\subsection{Logic Gates}

\textbf{Basic Gates:}
\begin{itemize}
    \item \textbf{AND:} Output 1 only if all inputs are 1
    \item \textbf{OR:} Output 1 if at least one input is 1
    \item \textbf{NOT:} Inverts the input
    \item \textbf{NAND:} NOT-AND (Universal gate)
    \item \textbf{NOR:} NOT-OR (Universal gate)
    \item \textbf{XOR:} Output 1 if inputs are different
    \item \textbf{XNOR:} Output 1 if inputs are same
\end{itemize}

\begin{notebox}
NAND and NOR are universal gates - any Boolean function can be implemented using only NAND or only NOR gates.
\end{notebox}

\subsection{Combinational Circuits}

\begin{defbox}
\textbf{Combinational Circuit:} Output depends only on current inputs (no memory).
\end{defbox}

\textbf{Key Combinational Circuits:}

\textbf{1. Multiplexer (MUX):} Selects one of many inputs based on select lines.
\begin{itemize}
    \item $2^n$-to-1 MUX requires n select lines
    \item Can implement any Boolean function
\end{itemize}

\textbf{2. Demultiplexer (DEMUX):} Routes one input to one of many outputs.

\textbf{3. Encoder:} Converts $2^n$ inputs to n-bit output.

\textbf{4. Decoder:} Converts n-bit input to $2^n$ outputs.

\textbf{5. Half Adder:}
\begin{itemize}
    \item Sum = $A \oplus B$
    \item Carry = $A \cdot B$
\end{itemize}

\textbf{6. Full Adder:}
\begin{itemize}
    \item Sum = $A \oplus B \oplus C_{in}$
    \item Carry = $AB + BC_{in} + AC_{in}$
\end{itemize}

\subsection{Sequential Circuits}

\begin{defbox}
\textbf{Sequential Circuit:} Output depends on current inputs and previous state (has memory).
\end{defbox}

\subsubsection{Flip-Flops}

\textbf{SR Flip-Flop:}
\begin{itemize}
    \item S=1, R=0: Set (Q=1)
    \item S=0, R=1: Reset (Q=0)
    \item S=0, R=0: No change
    \item S=1, R=1: Invalid state
\end{itemize}

\textbf{D Flip-Flop:}
\begin{itemize}
    \item Output Q = D on clock edge
    \item No invalid states
\end{itemize}

\textbf{JK Flip-Flop:}
\begin{itemize}
    \item J=0, K=0: No change
    \item J=0, K=1: Reset (Q=0)
    \item J=1, K=0: Set (Q=1)
    \item J=1, K=1: Toggle
\end{itemize}

\textbf{T Flip-Flop:}
\begin{itemize}
    \item T=0: No change
    \item T=1: Toggle
\end{itemize}

\subsubsection{Counters and Registers}

\textbf{Counter:} A sequential circuit that goes through a prescribed sequence of states.
\begin{itemize}
    \item \textbf{Asynchronous (Ripple):} Clock input cascaded through flip-flops
    \item \textbf{Synchronous:} All flip-flops clocked simultaneously
    \item \textbf{Modulo-N:} Counts from 0 to N-1
\end{itemize}

\textbf{Register:} A group of flip-flops used to store multiple bits.
\begin{itemize}
    \item \textbf{Shift Register:} Data shifted left or right
    \item \textbf{SISO:} Serial In Serial Out
    \item \textbf{SIPO:} Serial In Parallel Out
    \item \textbf{PISO:} Parallel In Serial Out
    \item \textbf{PIPO:} Parallel In Parallel Out
\end{itemize}

\subsection{Minimization Techniques}

\textbf{Karnaugh Map (K-Map):}
\begin{itemize}
    \item Graphical method for simplifying Boolean expressions
    \item Group adjacent 1s in powers of 2 (1, 2, 4, 8, ...)
    \item Larger groups lead to simpler expressions
\end{itemize}

\textbf{Quine-McCluskey Method:}
\begin{itemize}
    \item Tabular method for minimization
    \item Can handle any number of variables
    \item Finds all prime implicants
\end{itemize}

%==============================================================================
\section{Computer Organization and Architecture}

\subsection{Number Systems}

\textbf{Number Representation:}
\begin{itemize}
    \item \textbf{Binary:} Base 2 (0, 1)
    \item \textbf{Octal:} Base 8 (0-7)
    \item \textbf{Decimal:} Base 10 (0-9)
    \item \textbf{Hexadecimal:} Base 16 (0-9, A-F)
\end{itemize}

\textbf{Signed Number Representation:}
\begin{itemize}
    \item \textbf{Sign-Magnitude:} MSB for sign, rest for magnitude
    \item \textbf{1's Complement:} Invert all bits for negative
    \item \textbf{2's Complement:} 1's complement + 1 (most common)
\end{itemize}

\begin{notebox}
In n-bit 2's complement:
\begin{itemize}
    \item Range: $-2^{n-1}$ to $2^{n-1}-1$
    \item Single representation for zero
    \item Easy arithmetic operations
\end{itemize}
\end{notebox}

\textbf{Floating Point (IEEE 754):}
\[
    (-1)^{\text{sign}} \times (1 + \text{mantissa}) \times 2^{\text{exponent - bias}}
\]

\textbf{Single Precision (32-bit):}
\begin{itemize}
    \item 1 bit sign
    \item 8 bits exponent (bias = 127)
    \item 23 bits mantissa
\end{itemize}

\textbf{Double Precision (64-bit):}
\begin{itemize}
    \item 1 bit sign
    \item 11 bits exponent (bias = 1023)
    \item 52 bits mantissa
\end{itemize}

\subsection{Computer Architecture}

\subsubsection{Von Neumann Architecture}

\textbf{Components:}
\begin{enumerate}
    \item \textbf{CPU:} Central Processing Unit
    \item \textbf{Memory:} Stores instructions and data
    \item \textbf{I/O Devices:} Input/Output devices
    \item \textbf{Bus:} Communication pathway
\end{enumerate}

\textbf{Characteristics:}
\begin{itemize}
    \item Stored program concept
    \item Same memory for instructions and data
    \item Sequential instruction execution
    \item Von Neumann bottleneck: CPU and memory communicate through single bus
\end{itemize}

\subsubsection{Harvard Architecture}

\textbf{Characteristics:}
\begin{itemize}
    \item Separate memory for instructions and data
    \item Separate buses for instructions and data
    \item Can fetch instruction and data simultaneously
    \item Used in DSP processors
\end{itemize}

\subsection{CPU Organization}

\subsubsection{Registers}

\textbf{Types of Registers:}
\begin{itemize}
    \item \textbf{Program Counter (PC):} Holds address of next instruction
    \item \textbf{Instruction Register (IR):} Holds current instruction
    \item \textbf{Memory Address Register (MAR):} Holds memory address
    \item \textbf{Memory Data Register (MDR):} Holds data from/to memory
    \item \textbf{Accumulator (AC):} Temporary storage for ALU operations
    \item \textbf{Stack Pointer (SP):} Points to top of stack
    \item \textbf{General Purpose Registers:} Used for various operations
\end{itemize}

\subsubsection{Instruction Cycle}

\begin{defbox}
\textbf{Instruction Cycle Phases:}
\begin{enumerate}
    \item \textbf{Fetch:} Get instruction from memory
    \item \textbf{Decode:} Determine operation and operands
    \item \textbf{Execute:} Perform the operation
    \item \textbf{Store:} Write result back if needed
\end{enumerate}
\end{defbox}

\subsubsection{Addressing Modes}

\begin{enumerate}
    \item \textbf{Immediate:} Operand is part of instruction
    \item \textbf{Direct:} Address of operand in instruction
    \item \textbf{Indirect:} Address of address in instruction
    \item \textbf{Register:} Operand in register
    \item \textbf{Register Indirect:} Register contains address
    \item \textbf{Indexed:} Base address + index
    \item \textbf{Base Register:} Base register + offset
    \item \textbf{Relative:} PC + offset
\end{enumerate}

\subsubsection{Instruction Set Architecture}

\textbf{CISC (Complex Instruction Set Computer):}
\begin{itemize}
    \item Many specialized instructions
    \item Variable length instructions
    \item Multiple addressing modes
    \item Example: x86
\end{itemize}

\textbf{RISC (Reduced Instruction Set Computer):}
\begin{itemize}
    \item Simple, uniform instructions
    \item Fixed length instructions
    \item Load-store architecture
    \item Example: ARM, MIPS
\end{itemize}

\subsection{Pipeline}

\begin{defbox}
\textbf{Pipelining:} Technique to improve instruction throughput by overlapping execution of multiple instructions.
\end{defbox}

\textbf{Classic 5-Stage Pipeline:}
\begin{enumerate}
    \item \textbf{IF:} Instruction Fetch
    \item \textbf{ID:} Instruction Decode
    \item \textbf{EX:} Execute
    \item \textbf{MEM:} Memory Access
    \item \textbf{WB:} Write Back
\end{enumerate}

\begin{formulabox}
\textbf{Pipeline Performance:}
\begin{align*}
    \text{Speedup} &= \frac{\text{Non-pipelined time}}{\text{Pipelined time}} \\
    \text{Ideal Speedup} &= \text{Number of stages}
\end{align*}
\end{formulabox}

\textbf{Pipeline Hazards:}
\begin{enumerate}
    \item \textbf{Structural Hazard:} Resource conflict (e.g., single memory for data and instruction)
    \item \textbf{Data Hazard:} Instruction depends on result of previous instruction
    \begin{itemize}
        \item RAW (Read After Write) - True dependency
        \item WAR (Write After Read) - Anti dependency
        \item WAW (Write After Write) - Output dependency
    \end{itemize}
    \item \textbf{Control Hazard:} Branch instructions disrupt pipeline flow
\end{enumerate}

\textbf{Hazard Solutions:}
\begin{itemize}
    \item \textbf{Stalling:} Insert bubbles (NOPs) in pipeline
    \item \textbf{Forwarding (Bypassing):} Pass result directly to next stage
    \item \textbf{Branch Prediction:} Predict branch outcome
    \item \textbf{Delayed Branch:} Execute instruction after branch
\end{itemize}

\subsection{Memory Hierarchy}

\begin{defbox}
\textbf{Memory Hierarchy:} Organization of memory from fastest/smallest to slowest/largest.
\end{defbox}

\textbf{Hierarchy (top to bottom):}
\begin{enumerate}
    \item Registers (fastest, smallest)
    \item Cache (L1, L2, L3)
    \item Main Memory (RAM)
    \item Secondary Storage (HDD, SSD)
    \item Tertiary Storage (slowest, largest)
\end{enumerate}

\subsubsection{Cache Memory}

\begin{defbox}
\textbf{Cache:} Small, fast memory between CPU and main memory.
\end{defbox}

\textbf{Cache Performance Metrics:}
\begin{align*}
    \text{Hit Ratio} &= \frac{\text{Number of hits}}{\text{Total accesses}} \\
    \text{Miss Ratio} &= 1 - \text{Hit Ratio} \\
    \text{Average Access Time} &= \text{Hit Time} + \text{Miss Ratio} \times \text{Miss Penalty}
\end{align*}

\textbf{Cache Mapping Techniques:}

\textbf{1. Direct Mapped:}
\begin{itemize}
    \item Each block maps to exactly one cache line
    \item Cache line = (Block address) mod (Number of cache lines)
    \item Simple, fast, but high conflict misses
\end{itemize}

\textbf{2. Fully Associative:}
\begin{itemize}
    \item Block can go in any cache line
    \item Requires comparison with all tags
    \item Flexible but expensive
\end{itemize}

\textbf{3. Set Associative (n-way):}
\begin{itemize}
    \item Cache divided into sets
    \item Block maps to a set, can go in any line within set
    \item Compromise between direct and fully associative
\end{itemize}

\textbf{Replacement Policies:}
\begin{itemize}
    \item \textbf{LRU (Least Recently Used):} Replace block not used for longest time
    \item \textbf{FIFO (First In First Out):} Replace oldest block
    \item \textbf{Random:} Replace random block
    \item \textbf{LFU (Least Frequently Used):} Replace least frequently used block
\end{itemize}

\textbf{Write Policies:}
\begin{itemize}
    \item \textbf{Write-Through:} Write to both cache and memory
    \item \textbf{Write-Back:} Write only to cache, update memory on replacement
    \item \textbf{Write-Around:} Write directly to memory, bypass cache
\end{itemize}

\subsection{Virtual Memory}

\begin{defbox}
\textbf{Virtual Memory:} Technique that allows execution of programs larger than physical memory by using disk as extension of RAM.
\end{defbox}

\textbf{Paging:}
\begin{itemize}
    \item Divide memory into fixed-size pages
    \item Virtual address = Page number + Offset
    \item Page table maps virtual to physical addresses
\end{itemize}

\textbf{Translation Lookaside Buffer (TLB):}
\begin{itemize}
    \item Cache for page table entries
    \item Fast address translation
    \item Small, fully associative
\end{itemize}

\textbf{Page Replacement Algorithms:}
\begin{enumerate}
    \item \textbf{FIFO:} Replace oldest page
    \item \textbf{LRU:} Replace least recently used page
    \item \textbf{Optimal:} Replace page not used for longest time in future (theoretical)
    \item \textbf{Second Chance:} FIFO with reference bit
    \item \textbf{Clock:} Circular list with reference bit
\end{enumerate}

\begin{notebox}
\textbf{Belady's Anomaly:} In FIFO, increasing number of page frames can increase page faults.
\end{notebox}

\subsection{I/O Organization}

\textbf{I/O Methods:}
\begin{enumerate}
    \item \textbf{Programmed I/O:} CPU polls device status
    \item \textbf{Interrupt-Driven I/O:} Device interrupts CPU when ready
    \item \textbf{DMA (Direct Memory Access):} Device transfers data directly to/from memory
\end{enumerate}

\textbf{I/O Performance Metrics:}
\begin{itemize}
    \item \textbf{Throughput:} Amount of data transferred per unit time
    \item \textbf{Latency:} Time for single I/O operation
    \item \textbf{Bandwidth:} Maximum data transfer rate
\end{itemize}

%==============================================================================
\section{Programming and Data Structures}

\subsection{Programming Fundamentals}

\subsubsection{Time Complexity}

\begin{defbox}
\textbf{Big-O Notation:} $f(n) = O(g(n))$ if $\exists c, n_0$ such that $f(n) \leq c \cdot g(n)$ for all $n \geq n_0$.
\end{defbox}

\textbf{Common Time Complexities (from best to worst):}
\begin{itemize}
    \item $O(1)$ - Constant
    \item $O(\log n)$ - Logarithmic
    \item $O(n)$ - Linear
    \item $O(n \log n)$ - Linearithmic
    \item $O(n^2)$ - Quadratic
    \item $O(n^3)$ - Cubic
    \item $O(2^n)$ - Exponential
    \item $O(n!)$ - Factorial
\end{itemize}

\textbf{Other Notations:}
\begin{itemize}
    \item \textbf{Big-$\Omega$:} Lower bound
    \item \textbf{Big-$\Theta$:} Tight bound (both upper and lower)
    \item \textbf{Little-o:} Strict upper bound
\end{itemize}

\subsubsection{Recursion}

\begin{defbox}
\textbf{Recursion:} A function that calls itself directly or indirectly.
\end{defbox}

\textbf{Components of Recursion:}
\begin{enumerate}
    \item \textbf{Base Case:} Termination condition
    \item \textbf{Recursive Case:} Function calls itself with modified parameters
\end{enumerate}

\textbf{Types:}
\begin{itemize}
    \item \textbf{Direct Recursion:} Function calls itself
    \item \textbf{Indirect Recursion:} Function A calls B, B calls A
    \item \textbf{Tail Recursion:} Recursive call is last operation
\end{itemize}

\subsection{Arrays}

\begin{defbox}
\textbf{Array:} Collection of elements of same type stored in contiguous memory locations.
\end{defbox}

\textbf{Properties:}
\begin{itemize}
    \item Fixed size
    \item Random access: $O(1)$
    \item Address calculation: $\text{Base} + i \times \text{size of element}$
\end{itemize}

\textbf{2D Array Storage:}
\begin{itemize}
    \item \textbf{Row-Major:} $A[i][j] = \text{Base} + (i \times n + j) \times \text{size}$
    \item \textbf{Column-Major:} $A[i][j] = \text{Base} + (j \times m + i) \times \text{size}$
\end{itemize}

\subsection{Linked Lists}

\begin{defbox}
\textbf{Linked List:} Linear data structure where elements are stored in nodes, each containing data and pointer to next node.
\end{defbox}

\textbf{Types:}
\begin{enumerate}
    \item \textbf{Singly Linked List:} Each node points to next
    \item \textbf{Doubly Linked List:} Each node has pointers to next and previous
    \item \textbf{Circular Linked List:} Last node points to first
\end{enumerate}

\textbf{Operations Complexity:}
\begin{itemize}
    \item Access: $O(n)$
    \item Search: $O(n)$
    \item Insert at beginning: $O(1)$
    \item Insert at end: $O(n)$ (or $O(1)$ with tail pointer)
    \item Delete: $O(n)$
\end{itemize}

\subsection{Stacks}

\begin{defbox}
\textbf{Stack:} Linear data structure following LIFO (Last In First Out) principle.
\end{defbox}

\textbf{Operations:}
\begin{itemize}
    \item \textbf{Push:} Add element to top - $O(1)$
    \item \textbf{Pop:} Remove element from top - $O(1)$
    \item \textbf{Peek/Top:} View top element - $O(1)$
    \item \textbf{isEmpty:} Check if stack is empty - $O(1)$
\end{itemize}

\textbf{Applications:}
\begin{itemize}
    \item Function call management
    \item Expression evaluation and conversion
    \item Backtracking algorithms
    \item Undo mechanisms
\end{itemize}

\subsection{Queues}

\begin{defbox}
\textbf{Queue:} Linear data structure following FIFO (First In First Out) principle.
\end{defbox}

\textbf{Operations:}
\begin{itemize}
    \item \textbf{Enqueue:} Add element to rear - $O(1)$
    \item \textbf{Dequeue:} Remove element from front - $O(1)$
    \item \textbf{Front:} View front element - $O(1)$
    \item \textbf{isEmpty:} Check if queue is empty - $O(1)$
\end{itemize}

\textbf{Types:}
\begin{enumerate}
    \item \textbf{Simple Queue:} Basic FIFO
    \item \textbf{Circular Queue:} Rear wraps around to front
    \item \textbf{Priority Queue:} Elements have priorities
    \item \textbf{Deque (Double-ended Queue):} Insert/delete from both ends
\end{enumerate}

\subsection{Trees}

\begin{defbox}
\textbf{Tree:} Hierarchical data structure with nodes connected by edges. One node is root, others are children.
\end{defbox}

\textbf{Tree Terminology:}
\begin{itemize}
    \item \textbf{Root:} Top node
    \item \textbf{Parent:} Node with children
    \item \textbf{Child:} Node connected below
    \item \textbf{Leaf:} Node with no children
    \item \textbf{Height:} Longest path from root to leaf
    \item \textbf{Depth:} Distance from root
    \item \textbf{Degree:} Number of children
\end{itemize}

\subsubsection{Binary Tree}

\begin{defbox}
\textbf{Binary Tree:} Tree where each node has at most two children (left and right).
\end{defbox}

\textbf{Types:}
\begin{itemize}
    \item \textbf{Full Binary Tree:} Every node has 0 or 2 children
    \item \textbf{Complete Binary Tree:} All levels filled except possibly last, which is filled left to right
    \item \textbf{Perfect Binary Tree:} All internal nodes have 2 children, all leaves at same level
    \item \textbf{Balanced Binary Tree:} Height difference of left and right subtrees $\leq 1$
\end{itemize}

\textbf{Properties:}
\begin{itemize}
    \item Maximum nodes at level $l$: $2^l$
    \item Maximum nodes in tree of height $h$: $2^{h+1} - 1$
    \item Minimum height for $n$ nodes: $\log_2(n+1) - 1$
\end{itemize}

\textbf{Tree Traversals:}
\begin{enumerate}
    \item \textbf{Inorder (LNR):} Left, Node, Right
    \item \textbf{Preorder (NLR):} Node, Left, Right
    \item \textbf{Postorder (LRN):} Left, Right, Node
    \item \textbf{Level Order:} Level by level (uses queue)
\end{enumerate}

\subsubsection{Binary Search Tree (BST)}

\begin{defbox}
\textbf{BST:} Binary tree where for each node:
\begin{itemize}
    \item All values in left subtree $<$ node value
    \item All values in right subtree $>$ node value
\end{itemize}
\end{defbox}

\textbf{Operations:}
\begin{itemize}
    \item Search: $O(h)$ where h is height
    \item Insert: $O(h)$
    \item Delete: $O(h)$
    \item Best case: $h = \log n$ (balanced)
    \item Worst case: $h = n$ (skewed)
\end{itemize}

\begin{notebox}
Inorder traversal of BST gives sorted sequence.
\end{notebox}

\subsubsection{AVL Tree}

\begin{defbox}
\textbf{AVL Tree:} Self-balancing BST where height difference of left and right subtrees is at most 1.
\end{defbox}

\textbf{Balance Factor:} $BF = \text{Height(Left Subtree)} - \text{Height(Right Subtree)}$

\textbf{Valid values:} -1, 0, 1

\textbf{Rotations:}
\begin{enumerate}
    \item \textbf{Left Rotation (LL):} Right subtree is heavy
    \item \textbf{Right Rotation (RR):} Left subtree is heavy
    \item \textbf{Left-Right Rotation (LR):} Left subtree's right is heavy
    \item \textbf{Right-Left Rotation (RL):} Right subtree's left is heavy
\end{enumerate}

\textbf{Operations:} All operations in $O(\log n)$

\subsubsection{Red-Black Tree}

\begin{defbox}
\textbf{Red-Black Tree:} Self-balancing BST with following properties:
\begin{enumerate}
    \item Every node is red or black
    \item Root is black
    \item All leaves (NIL) are black
    \item Red node cannot have red children
    \item All paths from node to leaves have same number of black nodes
\end{enumerate}
\end{defbox}

\textbf{Operations:} All operations in $O(\log n)$

\textbf{Comparison with AVL:}
\begin{itemize}
    \item AVL is more rigidly balanced
    \item Red-Black allows slightly unbalanced trees
    \item Red-Black has faster insertion/deletion
    \item AVL has faster lookup
\end{itemize}

\subsubsection{B-Tree}

\begin{defbox}
\textbf{B-Tree of order m:} Self-balancing tree where:
\begin{itemize}
    \item Each node has at most m children
    \item Each non-leaf node (except root) has at least $\lceil m/2 \rceil$ children
    \item Root has at least 2 children (if not leaf)
    \item All leaves are at same level
\end{itemize}
\end{defbox}

\textbf{Used in:} Database indexing, file systems

\subsection{Heaps}

\begin{defbox}
\textbf{Heap:} Complete binary tree satisfying heap property:
\begin{itemize}
    \item \textbf{Max Heap:} Parent $\geq$ children
    \item \textbf{Min Heap:} Parent $\leq$ children
\end{itemize}
\end{defbox}

\textbf{Array Representation:}
\begin{itemize}
    \item Parent of node $i$: $(i-1)/2$
    \item Left child of node $i$: $2i + 1$
    \item Right child of node $i$: $2i + 2$
\end{itemize}

\textbf{Operations:}
\begin{itemize}
    \item Insert: $O(\log n)$ - Add at end, heapify up
    \item Delete Max/Min: $O(\log n)$ - Remove root, move last to root, heapify down
    \item Get Max/Min: $O(1)$
    \item Build Heap: $O(n)$
\end{itemize}

\textbf{Applications:}
\begin{itemize}
    \item Priority queues
    \item Heap sort
    \item Graph algorithms (Dijkstra, Prim)
\end{itemize}

\subsection{Hashing}

\begin{defbox}
\textbf{Hashing:} Technique to map data to fixed-size values using a hash function.
\end{defbox}

\textbf{Hash Function Properties:}
\begin{itemize}
    \item Deterministic
    \item Uniform distribution
    \item Fast computation
    \item Minimize collisions
\end{itemize}

\textbf{Common Hash Functions:}
\begin{enumerate}
    \item \textbf{Division Method:} $h(k) = k \mod m$
    \item \textbf{Multiplication Method:} $h(k) = \lfloor m(kA \mod 1) \rfloor$
    \item \textbf{Mid-Square Method:} Square key, take middle digits
    \item \textbf{Folding Method:} Divide key into parts, add them
\end{enumerate}

\textbf{Collision Resolution:}

\textbf{1. Open Addressing:}
\begin{itemize}
    \item \textbf{Linear Probing:} $h(k, i) = (h(k) + i) \mod m$
    \item \textbf{Quadratic Probing:} $h(k, i) = (h(k) + c_1i + c_2i^2) \mod m$
    \item \textbf{Double Hashing:} $h(k, i) = (h_1(k) + i \cdot h_2(k)) \mod m$
\end{itemize}

\textbf{2. Chaining:} Store colliding elements in linked list at same index

\textbf{Load Factor:} $\alpha = \frac{n}{m}$ where n = number of elements, m = table size

\begin{notebox}
Average case operations (with good hash function):
\begin{itemize}
    \item Insert: $O(1)$
    \item Search: $O(1)$
    \item Delete: $O(1)$
\end{itemize}
\end{notebox}

\subsection{Graphs}

\begin{defbox}
\textbf{Graph:} $G = (V, E)$ where V is set of vertices and E is set of edges.
\end{defbox}

\textbf{Types:}
\begin{itemize}
    \item \textbf{Directed/Undirected}
    \item \textbf{Weighted/Unweighted}
    \item \textbf{Cyclic/Acyclic}
    \item \textbf{Connected/Disconnected}
\end{itemize}

\textbf{Representation:}
\begin{enumerate}
    \item \textbf{Adjacency Matrix:} 2D array, $A[i][j] = 1$ if edge exists
    \begin{itemize}
        \item Space: $O(V^2)$
        \item Edge lookup: $O(1)$
        \item Finding neighbors: $O(V)$
    \end{itemize}
    \item \textbf{Adjacency List:} Array of lists
    \begin{itemize}
        \item Space: $O(V + E)$
        \item Edge lookup: $O(V)$ worst case
        \item Finding neighbors: $O(\text{degree})$
    \end{itemize}
\end{enumerate}

%==============================================================================
\section{Algorithms}

\subsection{Sorting Algorithms}

\subsubsection{Bubble Sort}
\begin{itemize}
    \item Repeatedly swap adjacent elements if in wrong order
    \item Time: $O(n^2)$ average and worst, $O(n)$ best
    \item Space: $O(1)$
    \item Stable
\end{itemize}

\subsubsection{Selection Sort}
\begin{itemize}
    \item Find minimum and place at beginning
    \item Time: $O(n^2)$ all cases
    \item Space: $O(1)$
    \item Not stable
\end{itemize}

\subsubsection{Insertion Sort}
\begin{itemize}
    \item Build sorted array one element at a time
    \item Time: $O(n^2)$ average and worst, $O(n)$ best
    \item Space: $O(1)$
    \item Stable
    \item Good for small arrays and nearly sorted data
\end{itemize}

\subsubsection{Merge Sort}
\begin{itemize}
    \item Divide array, sort halves, merge
    \item Time: $O(n \log n)$ all cases
    \item Space: $O(n)$
    \item Stable
    \item Divide and conquer
\end{itemize}

\subsubsection{Quick Sort}
\begin{itemize}
    \item Pick pivot, partition, recursively sort
    \item Time: $O(n \log n)$ average, $O(n^2)$ worst
    \item Space: $O(\log n)$ average (recursion stack)
    \item Not stable
    \item In-place
\end{itemize}

\subsubsection{Heap Sort}
\begin{itemize}
    \item Build max heap, repeatedly extract max
    \item Time: $O(n \log n)$ all cases
    \item Space: $O(1)$
    \item Not stable
    \item In-place
\end{itemize}

\subsubsection{Counting Sort}
\begin{itemize}
    \item Count occurrences of each value
    \item Time: $O(n + k)$ where k is range
    \item Space: $O(k)$
    \item Stable
    \item Non-comparison sort
    \item Works when range is not too large
\end{itemize}

\subsubsection{Radix Sort}
\begin{itemize}
    \item Sort digit by digit using stable sort
    \item Time: $O(d(n + k))$ where d is number of digits
    \item Space: $O(n + k)$
    \item Stable
    \item Non-comparison sort
\end{itemize}

\subsection{Searching Algorithms}

\subsubsection{Linear Search}
\begin{itemize}
    \item Check each element sequentially
    \item Time: $O(n)$
    \item Works on unsorted data
\end{itemize}

\subsubsection{Binary Search}
\begin{itemize}
    \item Divide sorted array in half repeatedly
    \item Time: $O(\log n)$
    \item Requires sorted array
\end{itemize}

\subsection{Graph Algorithms}

\subsubsection{Breadth-First Search (BFS)}
\begin{itemize}
    \item Explore level by level using queue
    \item Time: $O(V + E)$
    \item Space: $O(V)$
    \item Finds shortest path in unweighted graph
\end{itemize}

\subsubsection{Depth-First Search (DFS)}
\begin{itemize}
    \item Explore as far as possible before backtracking
    \item Time: $O(V + E)$
    \item Space: $O(V)$ (recursion stack)
    \item Used for topological sort, cycle detection
\end{itemize}

\subsubsection{Dijkstra's Algorithm}
\begin{itemize}
    \item Finds shortest path from source to all vertices
    \item Time: $O((V + E) \log V)$ with min-heap
    \item Does not work with negative weights
    \item Greedy algorithm
\end{itemize}

\subsubsection{Bellman-Ford Algorithm}
\begin{itemize}
    \item Finds shortest path from source
    \item Time: $O(VE)$
    \item Works with negative weights
    \item Detects negative cycles
    \item Dynamic programming
\end{itemize}

\subsubsection{Floyd-Warshall Algorithm}
\begin{itemize}
    \item All pairs shortest path
    \item Time: $O(V^3)$
    \item Space: $O(V^2)$
    \item Works with negative weights
    \item Dynamic programming
\end{itemize}

\subsubsection{Prim's Algorithm}
\begin{itemize}
    \item Finds Minimum Spanning Tree
    \item Time: $O((V + E) \log V)$ with min-heap
    \item Greedy algorithm
    \item Starts from single vertex
\end{itemize}

\subsubsection{Kruskal's Algorithm}
\begin{itemize}
    \item Finds Minimum Spanning Tree
    \item Time: $O(E \log E)$ or $O(E \log V)$
    \item Greedy algorithm
    \item Sorts edges by weight
    \item Uses Union-Find data structure
\end{itemize}

\subsubsection{Topological Sort}
\begin{itemize}
    \item Linear ordering of vertices in DAG
    \item Time: $O(V + E)$
    \item Uses DFS or Kahn's algorithm (BFS)
    \item Used in scheduling problems
\end{itemize}

\subsection{Dynamic Programming}

\begin{defbox}
\textbf{Dynamic Programming:} Solve complex problems by breaking down into simpler subproblems, storing results to avoid recomputation.
\end{defbox}

\textbf{Characteristics:}
\begin{enumerate}
    \item \textbf{Optimal Substructure:} Optimal solution contains optimal solutions to subproblems
    \item \textbf{Overlapping Subproblems:} Same subproblems solved multiple times
\end{enumerate}

\textbf{Approaches:}
\begin{itemize}
    \item \textbf{Top-Down (Memoization):} Recursion with caching
    \item \textbf{Bottom-Up (Tabulation):} Iterative, fill table
\end{itemize}

\textbf{Classic Problems:}
\begin{itemize}
    \item Fibonacci Numbers
    \item Longest Common Subsequence (LCS)
    \item Longest Increasing Subsequence (LIS)
    \item 0/1 Knapsack
    \item Matrix Chain Multiplication
    \item Edit Distance
    \item Coin Change
    \item Rod Cutting
\end{itemize}

\subsection{Greedy Algorithms}

\begin{defbox}
\textbf{Greedy Algorithm:} Make locally optimal choice at each step, hoping to find global optimum.
\end{defbox}

\textbf{When Greedy Works:}
\begin{enumerate}
    \item \textbf{Greedy Choice Property:} Local optimal choice leads to global optimal
    \item \textbf{Optimal Substructure:} Optimal solution contains optimal subproblem solutions
\end{enumerate}

\textbf{Classic Problems:}
\begin{itemize}
    \item Activity Selection
    \item Fractional Knapsack
    \item Huffman Coding
    \item Minimum Spanning Tree (Prim's, Kruskal's)
    \item Shortest Path (Dijkstra's)
\end{itemize}

\subsection{Divide and Conquer}

\begin{defbox}
\textbf{Divide and Conquer:} Break problem into smaller subproblems, solve recursively, combine solutions.
\end{defbox}

\textbf{Steps:}
\begin{enumerate}
    \item \textbf{Divide:} Break into subproblems
    \item \textbf{Conquer:} Solve subproblems recursively
    \item \textbf{Combine:} Merge solutions
\end{enumerate}

\textbf{Examples:}
\begin{itemize}
    \item Merge Sort
    \item Quick Sort
    \item Binary Search
    \item Strassen's Matrix Multiplication
\end{itemize}

\subsection{Backtracking}

\begin{defbox}
\textbf{Backtracking:} Build solution incrementally, abandon solution if it cannot lead to valid answer.
\end{defbox}

\textbf{Classic Problems:}
\begin{itemize}
    \item N-Queens
    \item Sudoku
    \item Graph Coloring
    \item Hamiltonian Path
    \item Subset Sum
\end{itemize}

\subsection{String Algorithms}

\subsubsection{Pattern Matching}

\textbf{Naive Algorithm:}
\begin{itemize}
    \item Check pattern at each position
    \item Time: $O(nm)$ where n = text length, m = pattern length
\end{itemize}

\textbf{KMP (Knuth-Morris-Pratt):}
\begin{itemize}
    \item Use failure function to avoid redundant comparisons
    \item Time: $O(n + m)$
    \item Preprocessing: Build LPS (Longest Proper Prefix Suffix) array
\end{itemize}

\textbf{Rabin-Karp:}
\begin{itemize}
    \item Use hashing to match pattern
    \item Average time: $O(n + m)$
    \item Worst time: $O(nm)$
\end{itemize}

%==============================================================================
\section{Theory of Computation}

\subsection{Automata Theory}

\subsubsection{Finite Automata}

\begin{defbox}
\textbf{Deterministic Finite Automaton (DFA):} $M = (Q, \Sigma, \delta, q_0, F)$ where:
\begin{itemize}
    \item Q: Finite set of states
    \item $\Sigma$: Input alphabet
    \item $\delta$: Transition function $Q \times \Sigma \rightarrow Q$
    \item $q_0$: Start state
    \item F: Set of accept states
\end{itemize}
\end{defbox}

\begin{defbox}
\textbf{Nondeterministic Finite Automaton (NFA):} Similar to DFA but:
\begin{itemize}
    \item $\delta: Q \times \Sigma_\epsilon \rightarrow P(Q)$ where $\Sigma_\epsilon = \Sigma \cup \{\epsilon\}$
    \item Can have $\epsilon$-transitions
    \item Multiple transitions for same input
\end{itemize}
\end{defbox}

\begin{thmbox}
\textbf{Equivalence:} For every NFA, there exists an equivalent DFA.
\end{thmbox}

\textbf{Regular Languages Properties:}
\begin{itemize}
    \item Closed under union, concatenation, Kleene star
    \item Closed under complement, intersection
    \item Pumping lemma for regular languages
\end{itemize}

\begin{thmbox}
\textbf{Pumping Lemma (Regular):} If L is regular, then $\exists p$ (pumping length) such that $\forall s \in L$ with $|s| \geq p$, s can be divided into $s = xyz$ where:
\begin{enumerate}
    \item $|xy| \leq p$
    \item $|y| > 0$
    \item $\forall i \geq 0, xy^iz \in L$
\end{enumerate}
\end{thmbox}

\subsubsection{Regular Expressions}

\begin{defbox}
\textbf{Regular Expression:} Pattern describing regular language using:
\begin{itemize}
    \item Concatenation
    \item Union (|)
    \item Kleene star (*)
\end{itemize}
\end{defbox}

\begin{thmbox}
A language is regular if and only if it can be described by a regular expression.
\end{thmbox}

\subsubsection{Context-Free Grammars}

\begin{defbox}
\textbf{Context-Free Grammar (CFG):} $G = (V, \Sigma, R, S)$ where:
\begin{itemize}
    \item V: Variables (non-terminals)
    \item $\Sigma$: Terminals
    \item R: Production rules ($ A \rightarrow \alpha$)
    \item S: Start symbol
\end{itemize}
\end{defbox}

\textbf{Chomsky Normal Form (CNF):}
\begin{itemize}
    \item $A \rightarrow BC$ or $A \rightarrow a$
    \item All productions in one of these forms
\end{itemize}

\textbf{Derivations:}
\begin{itemize}
    \item \textbf{Leftmost:} Replace leftmost variable first
    \item \textbf{Rightmost:} Replace rightmost variable first
\end{itemize}

\begin{defbox}
\textbf{Ambiguity:} A grammar is ambiguous if a string has multiple parse trees.
\end{defbox}

\subsubsection{Pushdown Automata (PDA)}

\begin{defbox}
\textbf{PDA:} Finite automaton with stack.
\end{defbox}

\begin{thmbox}
A language is context-free if and only if it is recognized by a PDA.
\end{thmbox}

\begin{thmbox}
\textbf{Pumping Lemma (Context-Free):} If L is context-free, then $\exists p$ such that $\forall s \in L$ with $|s| \geq p$, s can be divided into $s = uvxyz$ where:
\begin{enumerate}
    \item $|vy| > 0$
    \item $|vxy| \leq p$
    \item $\forall i \geq 0, uv^ixy^iz \in L$
\end{enumerate}
\end{thmbox}

\subsubsection{Turing Machines}

\begin{defbox}
\textbf{Turing Machine:} $M = (Q, \Sigma, \Gamma, \delta, q_0, q_{accept}, q_{reject})$ where:
\begin{itemize}
    \item Q: States
    \item $\Sigma$: Input alphabet
    \item $\Gamma$: Tape alphabet ($\Sigma \subset \Gamma$)
    \item $\delta: Q \times \Gamma \rightarrow Q \times \Gamma \times \{L, R\}$
    \item $q_0$: Start state
    \item $q_{accept}, q_{reject}$: Accept/reject states
\end{itemize}
\end{defbox}

\textbf{Variants:}
\begin{itemize}
    \item Multi-tape TM
    \item Non-deterministic TM
    \item All equivalent in power to standard TM
\end{itemize}

\subsection{Computability Theory}

\subsubsection{Decidability}

\begin{defbox}
\textbf{Decidable Language:} A language for which there exists a TM that halts on all inputs and accepts strings in the language.
\end{defbox}

\begin{defbox}
\textbf{Recognizable (Recursively Enumerable):} A language for which there exists a TM that accepts all strings in the language (may not halt on strings not in language).
\end{defbox}

\begin{thmbox}
Every decidable language is recognizable, but not vice versa.
\end{thmbox}

\textbf{Undecidable Problems:}
\begin{itemize}
    \item \textbf{Halting Problem:} Determine if TM halts on given input
    \item \textbf{$A_{TM}$:} Does TM M accept input w?
    \item \textbf{Rice's Theorem:} Any non-trivial property of TM languages is undecidable
\end{itemize}

\subsubsection{Reducibility}

\begin{defbox}
\textbf{Mapping Reduction:} Language A is mapping reducible to B ($A \leq_m B$) if there exists computable function $f$ such that:
\[
    w \in A \Leftrightarrow f(w) \in B
\]
\end{defbox}

\textbf{Properties:}
\begin{itemize}
    \item If $A \leq_m B$ and B is decidable, then A is decidable
    \item If $A \leq_m B$ and A is undecidable, then B is undecidable
\end{itemize}

\subsection{Complexity Theory}

\subsubsection{Time Complexity Classes}

\begin{defbox}
\textbf{P:} Set of languages decidable in polynomial time on deterministic TM.
\end{defbox}

\begin{defbox}
\textbf{NP:} Set of languages decidable in polynomial time on non-deterministic TM (or verifiable in polynomial time).
\end{defbox}

\begin{defbox}
\textbf{NP-Complete:} A language L is NP-complete if:
\begin{enumerate}
    \item $L \in NP$
    \item Every language in NP is polynomial-time reducible to L
\end{enumerate}
\end{defbox}

\begin{defbox}
\textbf{NP-Hard:} Problems at least as hard as hardest problems in NP (need not be in NP).
\end{defbox}

\begin{notebox}
\textbf{P vs NP:} Open problem. Does $P = NP$? Most believe $P \neq NP$.
\end{notebox}

\textbf{NP-Complete Problems:}
\begin{itemize}
    \item SAT (Boolean Satisfiability)
    \item 3-SAT
    \item Clique
    \item Vertex Cover
    \item Hamiltonian Path
    \item Traveling Salesman Problem (decision version)
    \item Subset Sum
\end{itemize}

\begin{thmbox}
\textbf{Cook-Levin Theorem:} SAT is NP-complete.
\end{thmbox}

\subsubsection{Space Complexity Classes}

\begin{defbox}
\textbf{PSPACE:} Problems solvable using polynomial space.
\end{defbox}

\begin{thmbox}
$P \subseteq NP \subseteq PSPACE$
\end{thmbox}

%==============================================================================
\section{Compiler Design}

\subsection{Phases of Compiler}

\begin{defbox}
\textbf{Compiler:} Program that translates source code to target code (usually machine code).
\end{defbox}

\textbf{Phases:}
\begin{enumerate}
    \item \textbf{Lexical Analysis:} Convert source to tokens
    \item \textbf{Syntax Analysis (Parsing):} Check grammar, create parse tree
    \item \textbf{Semantic Analysis:} Check meaning, type checking
    \item \textbf{Intermediate Code Generation:} Generate intermediate representation
    \item \textbf{Code Optimization:} Improve efficiency
    \item \textbf{Code Generation:} Generate target code
\end{enumerate}

\subsection{Lexical Analysis}

\begin{defbox}
\textbf{Lexical Analyzer (Lexer/Scanner):} Reads source code and produces tokens.
\end{defbox}

\textbf{Token:} Basic unit (keyword, identifier, operator, etc.)

\textbf{Implementation:}
\begin{itemize}
    \item Uses finite automata
    \item Regular expressions define tokens
    \item Tools: Lex, Flex
\end{itemize}

\subsection{Syntax Analysis}

\begin{defbox}
\textbf{Parser:} Checks if token sequence follows grammar rules.
\end{defbox}

\textbf{Parsing Techniques:}

\textbf{1. Top-Down Parsing:}
\begin{itemize}
    \item Start from start symbol, try to reach input
    \item \textbf{Recursive Descent:} Recursive functions for each non-terminal
    \item \textbf{LL(1):} Left-to-right scan, Leftmost derivation, 1 lookahead
    \begin{itemize}
        \item Uses parsing table
        \item No left recursion
        \item No ambiguity
    \end{itemize}
\end{itemize}

\textbf{2. Bottom-Up Parsing:}
\begin{itemize}
    \item Start from input, reduce to start symbol
    \item \textbf{LR Parsing:} Left-to-right scan, Rightmost derivation in reverse
    \item Types: SLR, LALR, CLR (increasing power)
    \item Uses parsing table and stack
    \item More powerful than LL
\end{itemize}

\textbf{Parse Tree vs Abstract Syntax Tree (AST):}
\begin{itemize}
    \item Parse tree: Shows all grammar rules
    \item AST: Simplified, shows structure without grammar details
\end{itemize}

\subsection{Semantic Analysis}

\textbf{Tasks:}
\begin{itemize}
    \item \textbf{Type Checking:} Ensure type consistency
    \item \textbf{Scope Resolution:} Check variable declarations
    \item \textbf{Symbol Table:} Store information about identifiers
\end{itemize}

\textbf{Symbol Table:}
\begin{itemize}
    \item Stores: name, type, scope, memory location
    \item Operations: insert, lookup, delete
    \item Implementation: hash table, tree
\end{itemize}

\subsection{Intermediate Code Generation}

\textbf{Common Forms:}
\begin{enumerate}
    \item \textbf{Three-Address Code:} $x = y \text{ op } z$
    \item \textbf{Abstract Syntax Tree}
    \item \textbf{Directed Acyclic Graph (DAG):} Identifies common subexpressions
\end{enumerate}

\subsection{Code Optimization}

\textbf{Types:}
\begin{itemize}
    \item \textbf{Machine Independent:}
    \begin{itemize}
        \item Constant folding
        \item Dead code elimination
        \item Common subexpression elimination
        \item Loop optimization (invariant code motion, strength reduction)
    \end{itemize}
    \item \textbf{Machine Dependent:}
    \begin{itemize}
        \item Register allocation
        \item Instruction scheduling
        \item Peephole optimization
    \end{itemize}
\end{itemize}

\subsection{Code Generation}

\textbf{Tasks:}
\begin{itemize}
    \item Instruction selection
    \item Register allocation
    \item Instruction ordering
\end{itemize}

\textbf{Register Allocation:}
\begin{itemize}
    \item Graph coloring approach
    \item Interference graph: nodes = variables, edge = live at same time
    \item k-colorable $\Rightarrow$ can use k registers
\end{itemize}

\subsection{Runtime Environments}

\textbf{Memory Organization:}
\begin{enumerate}
    \item \textbf{Code Area:} Stores instructions
    \item \textbf{Static Data:} Global variables
    \item \textbf{Stack:} Function calls, local variables
    \item \textbf{Heap:} Dynamic memory allocation
\end{enumerate}

\textbf{Activation Record (Stack Frame):}
\begin{itemize}
    \item Return address
    \item Parameters
    \item Local variables
    \item Saved registers
    \item Control link (previous frame pointer)
    \item Access link (for nested functions)
\end{itemize}

%==============================================================================
\section{Operating Systems}

\subsection{Process Management}

\subsubsection{Process}

\begin{defbox}
\textbf{Process:} Program in execution.
\end{defbox}

\textbf{Process States:}
\begin{enumerate}
    \item \textbf{New:} Being created
    \item \textbf{Ready:} Waiting for CPU
    \item \textbf{Running:} Executing
    \item \textbf{Waiting:} Waiting for I/O or event
    \item \textbf{Terminated:} Finished execution
\end{enumerate}

\textbf{Process Control Block (PCB):}
\begin{itemize}
    \item Process ID
    \item Process state
    \item Program counter
    \item CPU registers
    \item Memory management info
    \item I/O status
    \item Scheduling information
\end{itemize}

\subsubsection{CPU Scheduling}

\textbf{Scheduling Criteria:}
\begin{itemize}
    \item \textbf{CPU Utilization:} Keep CPU busy
    \item \textbf{Throughput:} Processes completed per unit time
    \item \textbf{Turnaround Time:} Time from submission to completion
    \item \textbf{Waiting Time:} Time spent in ready queue
    \item \textbf{Response Time:} Time from submission to first response
\end{itemize}

\textbf{Scheduling Algorithms:}

\textbf{1. First-Come First-Served (FCFS):}
\begin{itemize}
    \item Non-preemptive
    \item Simple but convoy effect
    \item Average waiting time often long
\end{itemize}

\textbf{2. Shortest Job First (SJF):}
\begin{itemize}
    \item Schedule process with smallest CPU burst
    \item Optimal for minimum average waiting time
    \item Can be preemptive (SRTF) or non-preemptive
    \item Problem: Starvation, predicting burst time
\end{itemize}

\textbf{3. Priority Scheduling:}
\begin{itemize}
    \item Each process has priority
    \item Problem: Starvation
    \item Solution: Aging (increase priority over time)
\end{itemize}

\textbf{4. Round Robin (RR):}
\begin{itemize}
    \item Preemptive, time quantum
    \item Fair, no starvation
    \item Performance depends on time quantum
\end{itemize}

\textbf{5. Multilevel Queue:}
\begin{itemize}
    \item Multiple queues with different priorities
    \item Each queue can have different algorithm
\end{itemize}

\textbf{6. Multilevel Feedback Queue:}
\begin{itemize}
    \item Processes can move between queues
    \item Aging implemented
    \item Most general and complex
\end{itemize}

\subsubsection{Process Synchronization}

\textbf{Critical Section Problem:}
\begin{itemize}
    \item \textbf{Mutual Exclusion:} Only one process in critical section
    \item \textbf{Progress:} Selection cannot be postponed indefinitely
    \item \textbf{Bounded Waiting:} Limit on waiting time
\end{itemize}

\textbf{Solutions:}

\textbf{1. Peterson's Solution:}
\begin{itemize}
    \item Software solution for 2 processes
    \item Uses turn and flag variables
\end{itemize}

\textbf{2. Semaphores:}
\begin{defbox}
\textbf{Semaphore:} Integer variable with two atomic operations:
\begin{itemize}
    \item \textbf{wait(S) or P(S):} while $S \leq 0$ do nothing; $S = S - 1$
    \item \textbf{signal(S) or V(S):} $S = S + 1$
\end{itemize}
\end{defbox}

\textbf{Types:}
\begin{itemize}
    \item \textbf{Binary Semaphore:} 0 or 1 (mutex)
    \item \textbf{Counting Semaphore:} Any integer
\end{itemize}

\textbf{3. Monitors:}
\begin{itemize}
    \item High-level synchronization construct
    \item Only one process active in monitor at a time
    \item Condition variables for waiting
\end{itemize}

\subsubsection{Deadlocks}

\begin{defbox}
\textbf{Deadlock:} Set of processes where each is waiting for resource held by another in the set.
\end{defbox}

\textbf{Necessary Conditions (all must hold):}
\begin{enumerate}
    \item \textbf{Mutual Exclusion:} Resources cannot be shared
    \item \textbf{Hold and Wait:} Process holds resources while waiting for more
    \item \textbf{No Preemption:} Resources cannot be forcibly taken
    \item \textbf{Circular Wait:} Circular chain of waiting processes
\end{enumerate}

\textbf{Handling Deadlocks:}

\textbf{1. Prevention:} Ensure at least one necessary condition cannot hold
\begin{itemize}
    \item Eliminate mutual exclusion (not always possible)
    \item Eliminate hold and wait (request all resources at once)
    \item Allow preemption
    \item Eliminate circular wait (order resources)
\end{itemize}

\textbf{2. Avoidance:} Use algorithm to avoid unsafe states
\begin{itemize}
    \item \textbf{Banker's Algorithm:}
    \begin{itemize}
        \item Keep system in safe state
        \item Safe state: sequence exists where all processes can complete
        \item Check if granting request keeps system safe
    \end{itemize}
\end{itemize}

\textbf{3. Detection and Recovery:}
\begin{itemize}
    \item Allow deadlocks
    \item Detect using resource allocation graph
    \item Recovery: Kill processes or preempt resources
\end{itemize}

\textbf{4. Ignorance:} Ostrich algorithm (used by many systems)

\subsection{Memory Management}

\subsubsection{Contiguous Allocation}

\textbf{Allocation Strategies:}
\begin{itemize}
    \item \textbf{First Fit:} Allocate first hole big enough
    \item \textbf{Best Fit:} Allocate smallest hole big enough
    \item \textbf{Worst Fit:} Allocate largest hole
\end{itemize}

\textbf{Fragmentation:}
\begin{itemize}
    \item \textbf{External:} Free memory divided into small blocks
    \item \textbf{Internal:} Allocated memory larger than requested
\end{itemize}

\subsubsection{Paging}

\begin{defbox}
\textbf{Paging:} Divide physical memory into fixed-size frames and logical memory into pages of same size.
\end{defbox}

\textbf{Address Translation:}
\[
    \text{Physical Address} = \text{Frame Number} \times \text{Page Size} + \text{Offset}
\]

\textbf{Page Table:}
\begin{itemize}
    \item Maps page number to frame number
    \item Each process has page table
    \item Stored in memory
\end{itemize}

\textbf{Translation Lookaside Buffer (TLB):}
\begin{itemize}
    \item Cache for page table
    \item Fast associative memory
    \item Contains recently used page-frame pairs
\end{itemize}

\begin{formulabox}
\textbf{Effective Access Time:}
\[
    EAT = \alpha \times (T + M) + (1-\alpha) \times (T + 2M)
\]
where $\alpha$ = TLB hit ratio, T = TLB access time, M = memory access time
\end{formulabox}

\subsubsection{Segmentation}

\begin{defbox}
\textbf{Segmentation:} Divide program into logical segments (code, data, stack, etc.).
\end{defbox}

\textbf{Segment Table:}
\begin{itemize}
    \item Base: Starting physical address
    \item Limit: Length of segment
\end{itemize}

\textbf{Paging vs Segmentation:}
\begin{itemize}
    \item Paging: Fixed size, invisible to user
    \item Segmentation: Variable size, logical division
\end{itemize}

\subsubsection{Virtual Memory}

\begin{defbox}
\textbf{Virtual Memory:} Technique allowing execution of processes not completely in memory.
\end{defbox}

\textbf{Benefits:}
\begin{itemize}
    \item Program larger than physical memory
    \item More processes in memory
    \item Less I/O for loading
\end{itemize}

\textbf{Demand Paging:}
\begin{itemize}
    \item Load page only when needed
    \item Valid/invalid bit in page table
    \item Page fault: required page not in memory
\end{itemize}

\textbf{Page Replacement Algorithms:} (Covered earlier in COA section)

\textbf{Thrashing:}
\begin{defbox}
\textbf{Thrashing:} Process spends more time paging than executing.
\end{defbox}

\textbf{Causes:}
\begin{itemize}
    \item Too many processes
    \item Insufficient memory
\end{itemize}

\textbf{Solutions:}
\begin{itemize}
    \item Working set model
    \item Page fault frequency
    \item Decrease degree of multiprogramming
\end{itemize}

\subsection{File Systems}

\subsubsection{File Concepts}

\begin{defbox}
\textbf{File:} Named collection of related information stored on secondary storage.
\end{defbox}

\textbf{File Attributes:}
\begin{itemize}
    \item Name, type, location
    \item Size, protection
    \item Time, date, user ID
\end{itemize}

\textbf{File Operations:}
\begin{itemize}
    \item Create, open, read, write
    \item Seek, close, delete
\end{itemize}

\subsubsection{Directory Structure}

\textbf{Types:}
\begin{enumerate}
    \item \textbf{Single-Level:} All files in one directory
    \item \textbf{Two-Level:} Separate directory for each user
    \item \textbf{Tree-Structured:} Hierarchical
    \item \textbf{Acyclic Graph:} Allows sharing (hard links)
    \item \textbf{General Graph:} Allows cycles (symbolic links)
\end{enumerate}

\subsubsection{File Allocation Methods}

\textbf{1. Contiguous Allocation:}
\begin{itemize}
    \item File occupies consecutive blocks
    \item Fast sequential access
    \item External fragmentation
\end{itemize}

\textbf{2. Linked Allocation:}
\begin{itemize}
    \item Each block points to next
    \item No external fragmentation
    \item Slow random access
    \item Example: FAT (File Allocation Table)
\end{itemize}

\textbf{3. Indexed Allocation:}
\begin{itemize}
    \item Index block contains pointers to data blocks
    \item Direct access
    \item Overhead of index block
    \item Example: UNIX inode
\end{itemize}

\subsubsection{Free Space Management}

\textbf{Methods:}
\begin{enumerate}
    \item \textbf{Bit Vector:} Each bit represents block (0=free, 1=allocated)
    \item \textbf{Linked List:} Link all free blocks
    \item \textbf{Grouping:} Store addresses of n free blocks in first free block
    \item \textbf{Counting:} Store starting address and count of contiguous free blocks
\end{enumerate}

\subsection{Disk Scheduling}

\textbf{Disk Access Time:}
\[
    \text{Access Time} = \text{Seek Time} + \text{Rotational Latency} + \text{Transfer Time}
\]

\textbf{Scheduling Algorithms:}

\textbf{1. FCFS:} Serve in order of arrival

\textbf{2. SSTF (Shortest Seek Time First):}
\begin{itemize}
    \item Service request closest to current head position
    \item May cause starvation
\end{itemize}

\textbf{3. SCAN (Elevator):}
\begin{itemize}
    \item Move head in one direction, service requests
    \item Reverse at end
\end{itemize}

\textbf{4. C-SCAN (Circular SCAN):}
\begin{itemize}
    \item Move in one direction only
    \item Jump to beginning at end
\end{itemize}

\textbf{5. LOOK and C-LOOK:}
\begin{itemize}
    \item Like SCAN/C-SCAN but reverse at last request, not end of disk
\end{itemize}

%==============================================================================
\section{Database Management Systems}

\subsection{Database Concepts}

\begin{defbox}
\textbf{Database:} Organized collection of interrelated data.
\end{defbox}

\begin{defbox}
\textbf{DBMS:} Software to create, maintain, and access databases.
\end{defbox}

\textbf{Advantages:}
\begin{itemize}
    \item Data independence
    \item Efficient data access
    \item Data integrity and security
    \item Concurrent access
    \item Crash recovery
\end{itemize}

\subsection{Database Models}

\textbf{1. Hierarchical:} Tree structure

\textbf{2. Network:} Graph structure

\textbf{3. Relational:} Tables (most common)

\textbf{4. Object-Oriented:} Objects with inheritance

\subsection{Relational Model}

\begin{defbox}
\textbf{Relation:} Table with rows (tuples) and columns (attributes).
\end{defbox}

\textbf{Keys:}
\begin{itemize}
    \item \textbf{Super Key:} Set of attributes uniquely identifying tuple
    \item \textbf{Candidate Key:} Minimal super key
    \item \textbf{Primary Key:} Chosen candidate key
    \item \textbf{Foreign Key:} References primary key of another relation
\end{itemize}

\textbf{Integrity Constraints:}
\begin{itemize}
    \item \textbf{Entity Integrity:} Primary key cannot be NULL
    \item \textbf{Referential Integrity:} Foreign key must reference existing primary key
    \item \textbf{Domain Constraint:} Attribute values from domain
\end{itemize}

\subsection{Relational Algebra}

\textbf{Unary Operations:}
\begin{itemize}
    \item \textbf{Selection ($\sigma$):} Select rows satisfying condition
    \item \textbf{Projection ($\pi$):} Select specific columns
    \item \textbf{Rename ($\rho$):} Rename relation/attributes
\end{itemize}

\textbf{Binary Operations:}
\begin{itemize}
    \item \textbf{Union ($\cup$):} Combine tuples from two relations
    \item \textbf{Intersection ($\cap$):} Common tuples
    \item \textbf{Difference ($-$):} Tuples in first but not second
    \item \textbf{Cartesian Product ($\times$):} All combinations
    \item \textbf{Join ($\bowtie$):} Combine relations based on condition
    \begin{itemize}
        \item Natural Join: Join on common attributes
        \item Theta Join: Join with condition
        \item Equi Join: Theta join with equality
        \item Outer Join: Include unmatched tuples (left, right, full)
    \end{itemize}
    \item \textbf{Division ($\div$):} Find tuples related to all tuples in other relation
\end{itemize}

\subsection{SQL}

\textbf{Data Definition Language (DDL):}
\begin{itemize}
    \item CREATE TABLE
    \item ALTER TABLE
    \item DROP TABLE
    \item CREATE INDEX
\end{itemize}

\textbf{Data Manipulation Language (DML):}
\begin{itemize}
    \item SELECT (query)
    \item INSERT
    \item UPDATE
    \item DELETE
\end{itemize}

\textbf{SELECT Statement:}
\begin{verbatim}
SELECT [DISTINCT] column_list
FROM table_list
WHERE condition
GROUP BY columns
HAVING group_condition
ORDER BY columns [ASC|DESC]
\end{verbatim}

\textbf{Aggregate Functions:}
\begin{itemize}
    \item COUNT, SUM, AVG, MAX, MIN
\end{itemize}

\textbf{Nested Queries:}
\begin{itemize}
    \item IN, NOT IN
    \item EXISTS, NOT EXISTS
    \item ANY, ALL
\end{itemize}

\textbf{Joins in SQL:}
\begin{itemize}
    \item INNER JOIN
    \item LEFT OUTER JOIN
    \item RIGHT OUTER JOIN
    \item FULL OUTER JOIN
    \item CROSS JOIN
\end{itemize}

\subsection{Normalization}

\begin{defbox}
\textbf{Normalization:} Process of organizing data to reduce redundancy and improve integrity.
\end{defbox}

\textbf{Functional Dependency:} $X \rightarrow Y$ means X determines Y

\textbf{Normal Forms:}

\textbf{1. First Normal Form (1NF):}
\begin{itemize}
    \item Atomic values (no multi-valued attributes)
    \item No repeating groups
\end{itemize}

\textbf{2. Second Normal Form (2NF):}
\begin{itemize}
    \item In 1NF
    \item No partial dependency (non-prime attribute depends on part of candidate key)
\end{itemize}

\textbf{3. Third Normal Form (3NF):}
\begin{itemize}
    \item In 2NF
    \item No transitive dependency (non-prime attribute depends on another non-prime attribute)
\end{itemize}

\textbf{Boyce-Codd Normal Form (BCNF):}
\begin{itemize}
    \item For every functional dependency $X \rightarrow Y$, X is a super key
    \item Stronger than 3NF
\end{itemize}

\textbf{4. Fourth Normal Form (4NF):}
\begin{itemize}
    \item In BCNF
    \item No multi-valued dependencies
\end{itemize}

\subsection{Transactions}

\begin{defbox}
\textbf{Transaction:} Logical unit of work (sequence of operations).
\end{defbox}

\textbf{ACID Properties:}
\begin{itemize}
    \item \textbf{Atomicity:} All or nothing
    \item \textbf{Consistency:} Database remains in consistent state
    \item \textbf{Isolation:} Concurrent transactions don't interfere
    \item \textbf{Durability:} Committed changes persist
\end{itemize}

\textbf{Transaction States:}
\begin{itemize}
    \item Active
    \item Partially Committed
    \item Committed
    \item Failed
    \item Aborted
\end{itemize}

\subsection{Concurrency Control}

\textbf{Problems:}
\begin{itemize}
    \item \textbf{Lost Update:} One transaction overwrites another's update
    \item \textbf{Dirty Read:} Read uncommitted data
    \item \textbf{Unrepeatable Read:} Different reads give different results
    \item \textbf{Phantom Read:} New rows appear in subsequent reads
\end{itemize}

\textbf{Serializability:}
\begin{defbox}
\textbf{Serial Schedule:} Transactions execute one after another.
\end{defbox}

\begin{defbox}
\textbf{Serializable Schedule:} Equivalent to some serial schedule.
\end{defbox}

\textbf{Conflict Serializability:}
\begin{itemize}
    \item Two operations conflict if from different transactions, on same item, at least one is write
    \item Schedule is conflict serializable if can be transformed to serial schedule by swapping non-conflicting operations
    \item Test: Create precedence graph, serializable if acyclic
\end{itemize}

\textbf{Locking Protocols:}

\textbf{1. Two-Phase Locking (2PL):}
\begin{itemize}
    \item \textbf{Growing Phase:} Acquire locks
    \item \textbf{Shrinking Phase:} Release locks
    \item Guarantees conflict serializability
    \item May cause deadlocks
\end{itemize}

\textbf{2. Strict 2PL:}
\begin{itemize}
    \item Hold all exclusive locks until commit
    \item Prevents cascading aborts
\end{itemize}

\textbf{3. Timestamp Ordering:}
\begin{itemize}
    \item Each transaction has timestamp
    \item Older transaction has priority
    \item No deadlocks but may have starvation
\end{itemize}

\subsection{Indexing}

\begin{defbox}
\textbf{Index:} Data structure to speed up data retrieval.
\end{defbox}

\textbf{Types:}

\textbf{1. Primary Index:}
\begin{itemize}
    \item On primary key
    \item Ordered file
\end{itemize}

\textbf{2. Clustering Index:}
\begin{itemize}
    \item On non-key field
    \item Records with same value clustered
\end{itemize}

\textbf{3. Secondary Index:}
\begin{itemize}
    \item On any field
    \item Dense index (entry for every record)
\end{itemize}

\textbf{B-Tree and B+-Tree:}
\begin{itemize}
    \item Balanced tree structures
    \item B+-Tree: All data in leaves, internal nodes only keys
    \item Efficient for range queries
    \item Most common in databases
\end{itemize}

\textbf{Hashing:}
\begin{itemize}
    \item Direct access using hash function
    \item Good for equality queries
    \item Not good for range queries
\end{itemize}

%==============================================================================
\section{Computer Networks}

\subsection{OSI Model}

\begin{defbox}
\textbf{OSI Model:} 7-layer reference model for network communication.
\end{defbox}

\textbf{Layers (top to bottom):}
\begin{enumerate}
    \item \textbf{Application:} User interface (HTTP, FTP, SMTP, DNS)
    \item \textbf{Presentation:} Data formatting, encryption (SSL/TLS)
    \item \textbf{Session:} Manages sessions between applications
    \item \textbf{Transport:} End-to-end communication (TCP, UDP)
    \item \textbf{Network:} Routing, logical addressing (IP)
    \item \textbf{Data Link:} Frame transmission, error detection (Ethernet, MAC)
    \item \textbf{Physical:} Physical transmission (cables, signals)
\end{enumerate}

\textbf{TCP/IP Model:}
\begin{enumerate}
    \item Application (combines OSI 5-7)
    \item Transport
    \item Internet (Network)
    \item Network Access (combines OSI 1-2)
\end{enumerate}

\subsection{Physical Layer}

\textbf{Transmission Media:}
\begin{itemize}
    \item \textbf{Guided:} Twisted pair, coaxial cable, fiber optic
    \item \textbf{Unguided:} Radio waves, microwaves, infrared
\end{itemize}

\textbf{Multiplexing:}
\begin{itemize}
    \item \textbf{FDM:} Frequency Division
    \item \textbf{TDM:} Time Division
    \item \textbf{WDM:} Wavelength Division
    \item \textbf{CDM:} Code Division
\end{itemize}

\subsection{Data Link Layer}

\textbf{Functions:}
\begin{itemize}
    \item Framing
    \item Physical addressing (MAC)
    \item Flow control
    \item Error control
    \item Access control
\end{itemize}

\textbf{Error Detection:}
\begin{itemize}
    \item \textbf{Parity Check:} Single bit parity
    \item \textbf{Checksum:} Sum of data segments
    \item \textbf{CRC (Cyclic Redundancy Check):} Polynomial division
\end{itemize}

\textbf{Flow Control:}
\begin{itemize}
    \item \textbf{Stop-and-Wait:} Send one frame, wait for ACK
    \item \textbf{Sliding Window:} Send multiple frames before ACK
\end{itemize}

\textbf{ARQ Protocols:}
\begin{itemize}
    \item \textbf{Stop-and-Wait ARQ}
    \item \textbf{Go-Back-N:} Retransmit all frames from error
    \item \textbf{Selective Repeat:} Retransmit only error frames
\end{itemize}

\textbf{Multiple Access Protocols:}
\begin{itemize}
    \item \textbf{ALOHA:} Pure and Slotted
    \item \textbf{CSMA:} Carrier Sense Multiple Access
    \item \textbf{CSMA/CD:} With Collision Detection (Ethernet)
    \item \textbf{CSMA/CA:} With Collision Avoidance (WiFi)
\end{itemize}

\textbf{Ethernet:}
\begin{itemize}
    \item MAC address: 48 bits
    \item Frame format: Preamble, destination, source, type, data, FCS
    \item CSMA/CD for access control
\end{itemize}

\subsection{Network Layer}

\textbf{Functions:}
\begin{itemize}
    \item Routing
    \item Logical addressing
    \item Packet forwarding
    \item Fragmentation
\end{itemize}

\subsubsection{IP Addressing}

\textbf{IPv4:}
\begin{itemize}
    \item 32-bit address
    \item Dotted decimal notation (e.g., 192.168.1.1)
    \item Classes A, B, C, D, E
\end{itemize}

\textbf{Classful Addressing:}
\begin{itemize}
    \item \textbf{Class A:} 0.0.0.0 to 127.255.255.255 (first bit 0)
    \item \textbf{Class B:} 128.0.0.0 to 191.255.255.255 (first bits 10)
    \item \textbf{Class C:} 192.0.0.0 to 223.255.255.255 (first bits 110)
    \item \textbf{Class D:} Multicast (first bits 1110)
    \item \textbf{Class E:} Reserved (first bits 1111)
\end{itemize}

\textbf{Subnetting:}
\begin{itemize}
    \item Divide network into subnetworks
    \item Subnet mask identifies network and host parts
    \item CIDR notation: IP/prefix length (e.g., 192.168.1.0/24)
\end{itemize}

\textbf{IPv6:}
\begin{itemize}
    \item 128-bit address
    \item Hexadecimal notation
    \item Larger address space
    \item Built-in security (IPSec)
\end{itemize}

\subsubsection{Routing Algorithms}

\textbf{Classification:}
\begin{itemize}
    \item \textbf{Static vs Dynamic}
    \item \textbf{Distance Vector vs Link State}
    \item \textbf{Interior vs Exterior}
\end{itemize}

\textbf{Distance Vector:}
\begin{itemize}
    \item Each router maintains distance table
    \item Periodic updates to neighbors
    \item Example: RIP (Routing Information Protocol)
    \item Uses Bellman-Ford algorithm
    \item Problem: Count-to-infinity
\end{itemize}

\textbf{Link State:}
\begin{itemize}
    \item Each router has complete network topology
    \item Flood link state information
    \item Example: OSPF (Open Shortest Path First)
    \item Uses Dijkstra's algorithm
\end{itemize}

\textbf{BGP (Border Gateway Protocol):}
\begin{itemize}
    \item Exterior gateway protocol
    \item Path vector protocol
    \item Routes between autonomous systems
\end{itemize}

\subsubsection{Network Layer Protocols}

\textbf{IP (Internet Protocol):}
\begin{itemize}
    \item Connectionless, unreliable
    \item Best effort delivery
    \item Fragmentation support
\end{itemize}

\textbf{ICMP (Internet Control Message Protocol):}
\begin{itemize}
    \item Error reporting and query
    \item Used by ping and traceroute
\end{itemize}

\textbf{ARP (Address Resolution Protocol):}
\begin{itemize}
    \item Maps IP to MAC address
    \item Broadcast ARP request
\end{itemize}

\textbf{DHCP (Dynamic Host Configuration Protocol):}
\begin{itemize}
    \item Automatically assign IP addresses
    \item Discover, Offer, Request, Acknowledge
\end{itemize}

\textbf{NAT (Network Address Translation):}
\begin{itemize}
    \item Translate private to public IP
    \item Conserve IP addresses
    \item Provides basic security
\end{itemize}

\subsection{Transport Layer}

\textbf{Functions:}
\begin{itemize}
    \item End-to-end communication
    \item Segmentation and reassembly
    \item Port addressing
    \item Connection control
    \item Flow and error control
\end{itemize}

\subsubsection{TCP (Transmission Control Protocol)}

\textbf{Characteristics:}
\begin{itemize}
    \item Connection-oriented
    \item Reliable
    \item Full-duplex
    \item Stream-oriented
    \item Flow and error control
\end{itemize}

\textbf{TCP Header:}
\begin{itemize}
    \item Source/Destination Port (16 bits each)
    \item Sequence Number (32 bits)
    \item Acknowledgment Number (32 bits)
    \item Flags: SYN, ACK, FIN, RST, PSH, URG
    \item Window Size
    \item Checksum
\end{itemize}

\textbf{Connection Establishment (3-Way Handshake):}
\begin{enumerate}
    \item Client: SYN
    \item Server: SYN + ACK
    \item Client: ACK
\end{enumerate}

\textbf{Connection Termination (4-Way):}
\begin{enumerate}
    \item Client: FIN
    \item Server: ACK
    \item Server: FIN
    \item Client: ACK
\end{enumerate}

\textbf{Flow Control:}
\begin{itemize}
    \item Sliding window protocol
    \item Receiver advertises window size
\end{itemize}

\textbf{Congestion Control:}
\begin{itemize}
    \item Slow Start
    \item Congestion Avoidance
    \item Fast Retransmit
    \item Fast Recovery
\end{itemize}

\subsubsection{UDP (User Datagram Protocol)}

\textbf{Characteristics:}
\begin{itemize}
    \item Connectionless
    \item Unreliable
    \item Lightweight
    \item Fast
    \item Message-oriented
\end{itemize}

\textbf{UDP Header (8 bytes):}
\begin{itemize}
    \item Source Port (16 bits)
    \item Destination Port (16 bits)
    \item Length (16 bits)
    \item Checksum (16 bits)
\end{itemize}

\textbf{Use Cases:}
\begin{itemize}
    \item DNS
    \item DHCP
    \item Streaming media
    \item Online gaming
    \item VoIP
\end{itemize}

\subsection{Application Layer}

\subsubsection{DNS (Domain Name System)}

\textbf{Function:} Translate domain names to IP addresses

\textbf{Hierarchy:}
\begin{itemize}
    \item Root servers
    \item Top-level domain (TLD) servers (.com, .org, etc.)
    \item Authoritative servers
\end{itemize}

\textbf{Query Types:}
\begin{itemize}
    \item Recursive: Server resolves completely
    \item Iterative: Server returns referral
\end{itemize}

\subsubsection{HTTP (Hypertext Transfer Protocol)}

\textbf{Characteristics:}
\begin{itemize}
    \item Stateless
    \item Application layer protocol
    \item Uses TCP (port 80)
\end{itemize}

\textbf{Methods:}
\begin{itemize}
    \item GET: Retrieve resource
    \item POST: Submit data
    \item PUT: Update resource
    \item DELETE: Delete resource
    \item HEAD: Get headers only
\end{itemize}

\textbf{Status Codes:}
\begin{itemize}
    \item 1xx: Informational
    \item 2xx: Success (200 OK)
    \item 3xx: Redirection (301 Moved, 304 Not Modified)
    \item 4xx: Client Error (400 Bad Request, 404 Not Found)
    \item 5xx: Server Error (500 Internal Server Error)
\end{itemize}

\textbf{HTTP/1.1 vs HTTP/2:}
\begin{itemize}
    \item HTTP/2: Binary protocol, multiplexing, server push
\end{itemize}

\textbf{HTTPS:}
\begin{itemize}
    \item HTTP over SSL/TLS
    \item Encrypted communication
    \item Port 443
\end{itemize}

\subsubsection{Email Protocols}

\textbf{SMTP (Simple Mail Transfer Protocol):}
\begin{itemize}
    \item Sending email
    \item Port 25
    \item Push protocol
\end{itemize}

\textbf{POP3 (Post Office Protocol):}
\begin{itemize}
    \item Download and delete from server
    \item Port 110
\end{itemize}

\textbf{IMAP (Internet Message Access Protocol):}
\begin{itemize}
    \item Access email on server
    \item Port 143
    \item Supports folders
\end{itemize}

\subsubsection{FTP (File Transfer Protocol)}

\textbf{Characteristics:}
\begin{itemize}
    \item Transfer files
    \item Uses two connections: control (port 21) and data (port 20)
    \item Active vs Passive mode
\end{itemize}

\subsection{Network Security}

\textbf{Cryptography:}
\begin{itemize}
    \item \textbf{Symmetric:} Same key for encryption and decryption (DES, AES)
    \item \textbf{Asymmetric:} Public and private keys (RSA, ECC)
\end{itemize}

\textbf{Digital Signatures:}
\begin{itemize}
    \item Ensure authenticity and non-repudiation
    \item Hash message, encrypt with private key
    \item Verify with public key
\end{itemize}

\textbf{Certificates:}
\begin{itemize}
    \item Issued by Certificate Authority (CA)
    \item Contains public key and identity
    \item X.509 standard
\end{itemize}

\textbf{Firewalls:}
\begin{itemize}
    \item Filter network traffic
    \item \textbf{Packet Filtering:} Based on headers
    \item \textbf{Stateful:} Track connections
    \item \textbf{Application Layer:} Deep packet inspection
\end{itemize}

\textbf{VPN (Virtual Private Network):}
\begin{itemize}
    \item Create secure tunnel over public network
    \item Encryption and authentication
\end{itemize}

%==============================================================================

\section{Tips for GATE Preparation}

\begin{notebox}
\textbf{Study Strategy:}
\begin{itemize}
    \item Understand concepts, don't just memorize
    \item Practice previous year questions
    \item Take mock tests regularly
    \item Focus on fundamentals first
    \item Make short notes for revision
    \item Solve problems from standard textbooks
    \item Time management is crucial
\end{itemize}
\end{notebox}

\begin{notebox}
\textbf{Subject-wise Priority:}
\begin{enumerate}
    \item High Priority: Data Structures, Algorithms, Operating Systems, DBMS, Computer Networks
    \item Medium Priority: Theory of Computation, Compiler Design, Computer Organization
    \item Lower Priority (but important): Digital Logic, Engineering Mathematics
\end{enumerate}
\end{notebox}

\begin{notebox}
\textbf{Important Formulas Summary:}
\begin{itemize}
    \item Time complexities of all algorithms
    \item Cache/TLB access time formulas
    \item Pipeline speedup calculations
    \item Page replacement algorithms
    \item Network layer calculations (IP, subnetting)
    \item Probability and combinatorics formulas
    \item Graph properties (Euler's, Handshaking lemma)
\end{itemize}
\end{notebox}

\textbf{Best wishes for your GATE preparation!}

\end{document}
